\chapter{OBJECTIVE}\label{obj}

The primary objectives of this laboratory exercise are as follows:

\begin{enumerate}
    \item \textbf{Translate Logical Design to Implementation Model:} To convert the conceptual and logical class diagrams into implementation-level specifications that include detailed attribute data types, method signatures with parameters and return types, and visibility indicators (public, private, protected).
    
    \item \textbf{Define Precise Class Specifications:} To establish complete class definitions for all domain entities in the Multi-Tenant E-Commerce System, ensuring each class contains all necessary attributes with appropriate data types (e.g., String, Integer, Float, Boolean, DateTime) and methods that represent the core operations of the system.
    
    \item \textbf{Model Relationships with Multiplicity:} To accurately represent the associations, aggregations, compositions, and inheritance relationships between classes, including precise multiplicity indicators (1, 0..*, 1..*, etc.) that define how objects interact with one another.
    
    \item \textbf{Ensure Multi-Tenant Data Isolation:} To design class structures that enforce the multi-tenant architecture requirement, ensuring that all tenant-specific data classes maintain proper references to the Tenant class, thereby enabling strict data segregation at the object level.
    
    \item \textbf{Establish a Blueprint for Implementation:} To create comprehensive class diagrams that serve as the definitive reference for developers during the coding phase, reducing ambiguity and ensuring consistency across all modules developed by different team members.
    
    \item \textbf{Facilitate Database Schema Design:} To provide a clear mapping from object-oriented classes to relational database tables, enabling efficient translation of the class model into the underlying database structure that supports both the central catalog database and isolated tenant databases.
\end{enumerate}
\clearpage
















